\section*{Chapitre \ref{ch:b2:ism-star-formation} --- Le milieu interstellaire et la formation des étoiles}

\begin{solution}{exo:b2:ism-star-formation:1}
$n = \num{2.7e-24}/\num{1.67e-24} = 1.6$ atomes d'hydrogène par centimètre cube. Dans l'air,
$n = P/k_{\mathrm B}T = 101325/(\num{1.381e-23}\times288) = \qty{2.55e25}{m^{-3}}$,
soit \num{2.5e19} molécules par centimètre cube~: environ $1.6\times10^{19}$ fois
plus.
\end{solution}

\begin{solution}{exo:b2:ism-star-formation:2}
$M_J \propto n^{-1/2}$~: diviser $n$ par 100 multiplie $M_J$ par 10, donc
$M_J = 54.6$ masses solaires à $10^2$ particules par centimètre cube. Le nuage contient $\num{24000}/54.6 \approx 440$ \omterm{def:b2:ism-star-formation:jeans}{masses de Jeans}~: il ne peut être stable dans son
ensemble et doit se
fragmenter.
\end{solution}

\begin{solution}{exo:b2:ism-star-formation:3}
$c_s = \sqrt{\num{1.381e-23}\times10/(2.33\times\num{1.661e-27})} =
\qty{189}{m.s^{-1}}$ à \qty{10}{K}~; $\sqrt{10}$ fois plus, soit \qty{597}{m.s^{-1}},
à \qty{100}{K}. Pour traverser \qty{0.1}{pc}~: $\num{3.086e15}/189 = \qty{1.63e13}{s}$,
soit \qty{0.52}{Ma}.
\end{solution}

\begin{solution}{exo:b2:ism-star-formation:4}
$\rho = 2.33\times\num{1.661e-27}\times n$ avec $n$ en \unit{m^{-3}}. Pour $10^5$
particules par centimètre cube, $\rho = \qty{3.87e-16}{kg.m^{-3}}$ et
$t_{\mathrm{ff}} = \sqrt{3\pi/(32G\rho)} = \qty{3.38e12}{s} = \qty{0.107}{Ma}$~; pour
$10^2$, $\rho$ est \num{1000} fois plus petit et $t_{\mathrm{ff}}$ est $\sqrt{1000} = 31.6$
fois plus long, \qty{3.38}{Ma}. Le \omterm{def:b2:ism-star-formation:molecular-cloud}{cœur dense} s'effondre 32 fois plus vite.
\end{solution}

\begin{solution}{exo:b2:ism-star-formation:5}
Avec $k_{\mathrm B}T/(\mu m_u) = c_s^2$, $(5k_{\mathrm B}T/G\mu m_u)^{3/2} =
5^{3/2}c_s^3G^{-3/2}$, ce qui donne la formule. Comme $M_J \propto \rho^{-1/2}$,
diviser $M_J$ par deux exige une densité quatre fois plus élevée. Comme $M_J \propto T^{3/2}$,
un réchauffement de \qty{10}{K} à \qty{20}{K} la multiplie par $2^{3/2} = 2.83$.
\end{solution}

\begin{solution}{exo:b2:ism-star-formation:6}
$t_{\mathrm{KH}} = GM^2/(RL) = \num{6.674e-11}\times(\num{1.989e30})^2/(2\times
\num{6.96e8}\times2\times\num{3.83e26}) = \qty{2.48e14}{s}$, soit
\qty{7.8}{Ma}. Elle est plus longue que la durée de vie de la plupart des disques (la moitié des disques ont disparu au bout de
environ \qty{3}{Ma})~: l'étoile se contracte encore quand son disque disparaît.
\end{solution}

\begin{solution}{exo:b2:ism-star-formation:7}
À \qty{10}{K}, $c_s^3/G = 189^3/\num{6.674e-11} = \qty{1.01e17}{kg.s^{-1}}$, soit
$\num{1.01e17}\times\num{3.156e7}/\num{1.989e30} = \num{1.6e-6}$ masses solaires par
an, et une masse solaire demande \qty{0.62}{Ma}. À \qty{20}{K}, le taux est
$2^{3/2} = 2.83$ fois plus élevé, soit \num{4.5e-6} masses solaires par an, et une
masse solaire demande \qty{0.22}{Ma}.
\end{solution}

\begin{solution}{exo:b2:ism-star-formation:8}
$r = \sqrt{L/(16\pi\sigma)}/T^2$ avec $\sqrt{\num{3.83e26}/(16\pi\times\num{5.67e-8})}
= \qty{1.159e16}{m.K^2}$~: à \qty{1300}{K}, $r = \qty{6.86e9}{m} = \qty{0.046}{au}$,
bien à l'intérieur de l'orbite de Mercure. À \qty{5.2}{au}, $T = 278/\sqrt{5.2} = \qty{122}{K}$.
\end{solution}

\begin{solution}{exo:b2:ism-star-formation:9}
$v = \sqrt{GM/r} = \sqrt{\num{6.674e-11}\times1.3\times\num{1.989e30}/\num{1.496e13}}
= \qty{3.40}{km.s^{-1}}$, et $P = 2\pi r/v = \qty{2.77e10}{s} = 877$ ans
(de manière équivalente $\sqrt{100^3/1.3}$ ans). En un million d'années, le gaz a effectué
$10^6/877 \approx \num{1140}$ tours~: le disque est vieux par rapport à sa propre rotation,
ce qui explique que ses anneaux soient des cercles réguliers.
\end{solution}

\begin{solution}{exo:b2:ism-star-formation:10}
$M = \int_0^R 4\pi r^2\rho_0R/r\,\mathrm dr = 2\pi\rho_0R^3$, et à l'intérieur de $r$,
$m(r) = 2\pi\rho_0Rr^2$. Avec $\mathrm dm = 4\pi\rho_0Rr\,\mathrm dr$,
\[
  U = -\int_0^R \frac{G\,m(r)}{r}\,\mathrm dm
    = -8\pi^2G\rho_0^2R^2\int_0^R r^2\,\mathrm dr = -\frac{8\pi^2}{3}G\rho_0^2R^5
    = -\frac23\,\frac{GM^2}{R},
\]
en utilisant $\rho_0 = M/(2\pi R^3)$. Elle vaut $\tfrac{2/3}{3/5} = 1.11$ fois la valeur uniforme~:
le nuage concentré est plus fortement lié et a besoin de \qty{11}{\percent} d'énergie thermique
supplémentaire pour être maintenu~; il s'effondre plus facilement.
\end{solution}

\begin{solution}{exo:b2:ism-star-formation:11}
Période $2\pi/\Omega = \qty{6.28e14}{s}$, soit environ \qty{20}{Ma}. Avec
$R = \qty{3.086e15}{m}$, $j = \Omega R^2 = \qty{9.52e16}{m^2.s^{-1}}$, et
$r_c = j^2/(GM) = \num{9.07e33}/\num{1.327e20} = \qty{6.83e13}{m} = \qty{457}{au}$.
En rotation solide, $j = \Omega r^2$ croît comme le carré de la distance à l'axe~:
le gaz venant de l'intérieur du cœur a moins de moment cinétique et tombe plus près de l'étoile.
\end{solution}

\begin{solution}{exo:b2:ism-star-formation:12}
$\bar t = 2.46$, $\bar f = 0.588$; $\sum(t - \bar t)^2 = 12.612$,
$\sum(t - \bar t)(f - \bar f) = -2.018$, donc $b = 2.018/12.612 = 0.160$ par million
d'années et $a = 0.588 + 0.160\times2.46 = 0.982$. La droite atteint $0.5$ à
$(0.982 - 0.5)/0.160 = \qty{3.0}{Ma}$ et zéro à $0.982/0.160 = \qty{6.1}{Ma}$, comme sur
la figure. Une fraction ne peut descendre sous zéro~: les amas plus âgés, comme NGC 1960 à
\qty{30}{Ma}, conservent une petite fraction voisine de zéro qu'une droite ne peut pas décrire.
\end{solution}

\begin{solution}{pb:b2:ism-star-formation:1}
\textbf{1.} $R = 0.08\times\num{3.086e16} = \qty{2.469e15}{m}$.

\textbf{2.} $N = \qty{1.5e26}{m^{-2}}$, donc $2.8\times\num{1.661e-27}\times\num{1.5e26} =
\qty{0.698}{kg.m^{-2}}$.

\textbf{3.} $M = 0.698\times\pi R^2 = 0.698\times\num{1.915e31} = \qty{1.34e31}{kg}$,
soit $6.7$ masses solaires.

\textbf{4.} $V = \tfrac43\pi R^3 = \qty{6.30e46}{m^3}$ et
$\rho = M/V = \qty{2.12e-16}{kg.m^{-3}}$.

\textbf{5.} $n = \rho/(2.33m_u) = \qty{5.48e10}{m^{-3}}$, soit \num{5.5e4} particules
par centimètre cube.

\textbf{6.} Oui~: sa densité se situe entre $10^4$ et $10^5$, son rayon entre 0.03 et
\qty{0.2}{pc}~; sa masse, 6.7 masses solaires, se situe au sommet de la fourchette des cœurs.

\textbf{7.} $M_J = 5.46\times(10^4/\num{5.48e4})^{1/2} = 2.33$ masses solaires.

\textbf{8.} $6.7/2.33 = 2.9$ \omterm{def:b2:ism-star-formation:jeans}{masses de Jeans}~: le cœur est instable et s'effondre.

\textbf{9.} $R_J = 0.1316\times(10^4/\num{5.48e4})^{1/2} = \qty{0.056}{pc}$, plus petit que
le rayon, \qty{0.08}{pc}~: même conclusion.

\textbf{10.} $2^{3/2}\times2.33 = 6.6$ masses solaires~: environ la masse du cœur. À
\qty{20}{K}, le cœur serait limite~; c'est le froid qui le fait s'effondrer.

\textbf{11.} $t_{\mathrm{ff}} = \sqrt{3\pi/(32\times\num{6.674e-11}\times\num{2.12e-16})}
= \qty{4.56e12}{s}$, soit environ \num{145000} ans.

\textbf{12.} $R/c_s = \num{2.469e15}/189 = \qty{1.31e13}{s}$, \num{414000} ans~: trois
fois plus long que la chute libre. Une onde de pression ne peut traverser le cœur avant
de s'être effondré~: la pression ne peut pas réagir assez vite pour arrêter l'effondrement.

\textbf{13.} $c_s = \qty{189}{m.s^{-1}}$ (\cref{exo:b2:ism-star-formation:3}).

\textbf{14.} $c_s^3/G = \qty{1.01e17}{kg.s^{-1}}$, \num{1.6e-6} masses solaires par an.

\textbf{15.} $1/\num{1.6e-6} = \num{624000}$ ans.

\textbf{16.} Un tiers de 6.7 vaut $2.2$~: deux étoiles d'une masse solaire.

\textbf{17.} $t_{\mathrm{KH}} = GM_\odot^2/(3R_\odot\times4L_\odot) = \qty{8.25e13}{s}$,
\qty{2.6}{Ma}.

\textbf{18.} Plus long, quatre fois le temps d'accrétion~: quand l'accrétion s'arrête, l'étoile est
encore en contraction, elle brille sans fusion~; c'est une \omterm{def:b2:ism-star-formation:t-tauri}{étoile T Tauri}, pas encore sur la séquence
principale.

\textbf{19.} $\Omega = \qty{0.4e3}{m.s^{-1}}/\qty{3.086e16}{m} = \qty{1.30e-14}{rad.s^{-1}}$.

\textbf{20.} $j = 1.30\times10^{-14}\times(\num{2.469e15})^2 = \qty{7.90e16}{m^2.s^{-1}}$.

\textbf{21.} $r_c = j^2/(GM_\odot) = \num{6.24e33}/\num{1.327e20} = \qty{4.70e13}{m}$, soit
\qty{314}{au}~: plus grand que le disque de HL Tauri, comme attendu pour du gaz venant du bord du
cœur.

\textbf{22.} $T = (L/16\pi\sigma r^2)^{1/4} = \qty{278}{K}$ à \qty{1}{au}, et
$278/\sqrt{30} = \qty{51}{K}$ à \qty{30}{au}.

\textbf{23.} $r = \num{1.159e16}/180^2 = \qty{3.58e11}{m}$, soit \qty{2.4}{au}, comme dans le texte.

\textbf{24.} Une étoile d'une masse solaire met
\textbf{\qty[text-series-to-math=true]{624}{} milliers d'années} pour se former, environ quatre
fois le \omterm{def:b2:ism-star-formation:free-fall}{temps de chute libre} de son cœur (\num{145000} ans)~: le gaz extérieur n'atteint
la \omterm{def:b2:ism-star-formation:protostar}{protoétoile} que progressivement, de l'intérieur vers l'extérieur.
\end{solution}
